\documentclass[10pt,a4paper]{amsart}
\usepackage{amsmath,amssymb,amsthm,graphicx}
% Retained minimal preamble for the Stacks Project chapter varieties.tex.
% Verbatim excerpt of the chapter preamble: Xy-pic setup (native preamble lines 14-18), the native theorem
% declarations with the mathematical macros (native preamble lines 154-217), and the project own no-op
% definitions of the reference/slogan/history display environments (native wrapper.tex lines 1-3) that the
% counted statement uses for its slogan. Copied mechanically; nothing added, removed or redefined.
% For commutative diagrams we use Xy-pic
\usepackage[all]{xy}

% We use 2cell for 2-commutative diagrams.
\xyoption{2cell}
% Theorem environments.
%
\theoremstyle{plain}
\newtheorem{theorem}[subsection]{Theorem}
\newtheorem{proposition}[subsection]{Proposition}
\newtheorem{lemma}[subsection]{Lemma}

\theoremstyle{definition}
\newtheorem{definition}[subsection]{Definition}
\newtheorem{example}[subsection]{Example}
\newtheorem{exercise}[subsection]{Exercise}
\newtheorem{situation}[subsection]{Situation}

\theoremstyle{remark}
\newtheorem{remark}[subsection]{Remark}
\newtheorem{remarks}[subsection]{Remarks}

\numberwithin{equation}{subsection}

% Macros
%
\def\lim{\mathop{\mathrm{lim}}\nolimits}
\def\colim{\mathop{\mathrm{colim}}\nolimits}
\def\Spec{\mathop{\mathrm{Spec}}}
\def\Hom{\mathop{\mathrm{Hom}}\nolimits}
\def\Ext{\mathop{\mathrm{Ext}}\nolimits}
\def\SheafHom{\mathop{\mathcal{H}\!\mathit{om}}\nolimits}
\def\SheafExt{\mathop{\mathcal{E}\!\mathit{xt}}\nolimits}
\def\Sch{\mathit{Sch}}
\def\Mor{\mathop{\mathrm{Mor}}\nolimits}
\def\Ob{\mathop{\mathrm{Ob}}\nolimits}
\def\Sh{\mathop{\mathit{Sh}}\nolimits}
\def\NL{\mathop{N\!L}\nolimits}
\def\CH{\mathop{\mathrm{CH}}\nolimits}
\def\proetale{{pro\text{-}\acute{e}tale}}
\def\etale{{\acute{e}tale}}
\def\QCoh{\mathit{QCoh}}
\def\Ker{\mathop{\mathrm{Ker}}}
\def\Im{\mathop{\mathrm{Im}}}
\def\Coker{\mathop{\mathrm{Coker}}}
\def\Coim{\mathop{\mathrm{Coim}}}

% Boxtimes
%
\DeclareMathSymbol{\boxtimes}{\mathbin}{AMSa}{"02}

%
% Macros for moduli stacks/spaces
%
\def\QCohstack{\mathcal{QC}\!\mathit{oh}}
\def\Cohstack{\mathcal{C}\!\mathit{oh}}
\def\Spacesstack{\mathcal{S}\!\mathit{paces}}
\def\Quotfunctor{\mathrm{Quot}}
\def\Hilbfunctor{\mathrm{Hilb}}
\def\Curvesstack{\mathcal{C}\!\mathit{urves}}
\def\Polarizedstack{\mathcal{P}\!\mathit{olarized}}
\def\Complexesstack{\mathcal{C}\!\mathit{omplexes}}
% \Pic is the operator that assigns to X its picard group, usage \Pic(X)
% \Picardstack_{X/B} denotes the Picard stack of X over B
% \Picardfunctor_{X/B} denotes the Picard functor of X over B
\def\Pic{\mathop{\mathrm{Pic}}\nolimits}
\def\Picardstack{\mathcal{P}\!\mathit{ic}}
\def\Picardfunctor{\mathrm{Pic}}
\def\Deformationcategory{\mathcal{D}\!\mathit{ef}}
\newenvironment{reference}{}{}
\newenvironment{slogan}{}{}
\newenvironment{history}{}{}
\usepackage{cleveref}
\crefname{originalreference}{original reference}{original references}
\Crefname{originalreference}{Original reference}{Original references}
\makeatletter
\newlabel{schemes-lemma-morphism-into-affine}{{01I1}{source}{}{sourcepointer-1}{}}
\newlabel{schemes-lemma-morphism-into-affine@cref}{{[originalreference][0][]01I1}{[1][0][]source}{}{}{}}
\makeatother
\crefname{originalreference}{original reference}{original references}
\Crefname{originalreference}{Original reference}{Original references}
\makeatletter
\newlabel{algebra-lemma-integral-no-inclusion}{{00GT}{source}{}{sourcepointer-2}{}}
\newlabel{algebra-lemma-integral-no-inclusion@cref}{{[originalreference][0][]00GT}{[1][0][]source}{}{}{}}
\makeatother
\crefname{originalreference}{original reference}{original references}
\Crefname{originalreference}{Original reference}{Original references}
\makeatletter
\newlabel{algebra-lemma-ring-with-only-minimal-primes}{{04MG}{source}{}{sourcepointer-3}{}}
\newlabel{algebra-lemma-ring-with-only-minimal-primes@cref}{{[originalreference][0][]04MG}{[1][0][]source}{}{}{}}
\makeatother
\crefname{originalreference}{original reference}{original references}
\Crefname{originalreference}{Original reference}{Original references}
\makeatletter
\newlabel{algebra-lemma-characterize-local-ring}{{00E9}{source}{}{sourcepointer-4}{}}
\newlabel{algebra-lemma-characterize-local-ring@cref}{{[originalreference][0][]00E9}{[1][0][]source}{}{}{}}
\makeatother
\crefname{originalreference}{original reference}{original references}
\Crefname{originalreference}{Original reference}{Original references}
\makeatletter
\newlabel{algebra-lemma-minimal-prime-reduced-ring}{{00EU}{source}{}{sourcepointer-5}{}}
\newlabel{algebra-lemma-minimal-prime-reduced-ring@cref}{{[originalreference][0][]00EU}{[1][0][]source}{}{}{}}
\makeatother
\title{The integral closure of the ground field in the global sections of a connected reduced scheme is a field}
\author{The Stacks Project Authors}
\date{stacks-project:varieties.tex:local-snapshot-20261011; GNU FDL 1.2 or later}
\begin{document}
\maketitle
\noindent\textit{Editorial note.} All delivered mathematical text is the author's own, copied verbatim from the pinned Stacks Project chapter file varieties.tex; no statement or proof step was rewritten, reordered or removed.\par\medskip
\noindent\textit{Editorial note.} The counted span cites five cross-chapter labels, each recorded in assembly.external\_references with its canonical tag as a navigation pointer.\par\medskip
\noindent\textit{Editorial note.} The proof cites no bibliography item inside the counted span.\par\medskip
\noindent\textit{Allowed standard prerequisites (editorial).}\begin{itemize}
\item Integral ring extensions have no inclusion of prime ideals and their spectra are totally disconnected. Stacks Project, Algebra, Lemmas algebra-lemma-integral-no-inclusion (tag 00GS) and algebra-lemma-ring-with-only-minimal-primes (tag 00G3). These give that every prime of the integral closure is minimal and maximal, which is what makes the constancy argument work.
\item Morphisms into affine schemes and the characterization of local rings. Stacks Project, Schemes, Lemma schemes-lemma-morphism-into-affine (tag 01I1); Algebra, Lemma algebra-lemma-characterize-local-ring (tag 00ER). The first factors the structure morphism through Spec(k') and the second identifies k' as local once all elements outside p are units.
\end{itemize}\medskip
\ifdefined\hypertarget\hypertarget{sourcepointer-1}{}\fi
\noindent\textit{Source reference.} Original reference 01I1: sources/stacks-chow/tags.txt, source lines 1962--1962. Cross-chapter standard prerequisite cited by the author inside the counted span; it is used as a pointer and its own proof is not part of this item.\par\medskip
\ifdefined\hypertarget\hypertarget{sourcepointer-2}{}\fi
\noindent\textit{Source reference.} Original reference 00GT: sources/stacks-chow/tags.txt, source lines 619--619. Cross-chapter standard prerequisite cited by the author inside the counted span; it is used as a pointer and its own proof is not part of this item.\par\medskip
\ifdefined\hypertarget\hypertarget{sourcepointer-3}{}\fi
\noindent\textit{Source reference.} Original reference 04MG: sources/stacks-chow/tags.txt, source lines 6042--6042. Cross-chapter standard prerequisite cited by the author inside the counted span; it is used as a pointer and its own proof is not part of this item.\par\medskip
\ifdefined\hypertarget\hypertarget{sourcepointer-4}{}\fi
\noindent\textit{Source reference.} Original reference 00E9: sources/stacks-chow/tags.txt, source lines 526--526. Cross-chapter standard prerequisite cited by the author inside the counted span; it is used as a pointer and its own proof is not part of this item.\par\medskip
\ifdefined\hypertarget\hypertarget{sourcepointer-5}{}\fi
\noindent\textit{Source reference.} Original reference 00EU: sources/stacks-chow/tags.txt, source lines 548--548. Cross-chapter standard prerequisite cited by the author inside the counted span; it is used as a pointer and its own proof is not part of this item.\par\medskip
\begin{lemma}
\label{lemma-integral-closure-ground-field}
Let $k$ be a field.
Let $X$ be a scheme over $k$ which is connected and reduced.
Then the integral closure of $k$ in $\Gamma(X, \mathcal{O}_X)$
is a field.
\end{lemma}

\begin{proof}
Let $k' \subset \Gamma(X, \mathcal{O}_X)$ be the integral closure of
$k$. Then $X \to \Spec(k)$ factors through $\Spec(k')$, see
Schemes, Lemma \ref{schemes-lemma-morphism-into-affine}.
As $X$ is reduced we see that $k'$ has no nonzero nilpotent elements.
As $k \to k'$ is integral we see that every prime ideal of $k'$ is
both a maximal ideal and a minimal prime,
and $\Spec(k')$ is totally disconnected, see
Algebra, Lemmas \ref{algebra-lemma-integral-no-inclusion} and
\ref{algebra-lemma-ring-with-only-minimal-primes}.
As $X$ is connected the morphism $X \to \Spec(k')$ is
constant, say with image the point corresponding to
$\mathfrak p \subset k'$. Then any $f \in k'$, $f \not \in \mathfrak p$
maps to an invertible element of $\mathcal{O}_X$. By definition
of $k'$ this then forces $f$ to be a unit of $k'$. Hence we see
that $k'$ is local with maximal ideal $\mathfrak p$, see
Algebra, Lemma \ref{algebra-lemma-characterize-local-ring}.
Since we've already seen that $k'$ is reduced this implies that
$k'$ is a field, see
Algebra, Lemma \ref{algebra-lemma-minimal-prime-reduced-ring}.
\end{proof}

\end{document}
